\SourceLeaf{194}mit diesem Gegenstande beschäftigt haben. Bei mir ist es über 30 Jahr, und ich glaube nicht,
dass jemand sich eben mit diesem $2^{\mathrm n}.$ Theil mehr beschäftigt haben könne als ich,
obgleich ich niemals darüber etwas bekannt gemacht habe. Die Annahme, dass die Summe der 3
Winkel kleiner sei als $180^{0}$, führt auf eine eigene, von der unsrigen (Euklidischen) ganz
verschiedene Geometrie, die in sich selbst durchaus consequent ist, und die ich für mich selbst
ganz befriedigend ausgebildet habe, so dass ich jede Aufgabe in derselben auflösen kann mit
Ausnahme der Bestimmung einer Constante, die sich a priori nicht ausmitteln lässt. Je grösser
man diese Constante annimmt, desto mehr nähert man sich der Euklidischen Geometrie und ein
unendlich grosser Werth macht beide zusammenfallen. Die Sätze jener Geometrie scheinen zum
Theil paradox, und dem Ungeübten ungereimt; bei genauerer ruhiger Überlegung findet man aber,
dass sie an sich durchaus nichts unmögliches enthalten. So z. B. können die drei Winkel eines
Dreiecks so klein werden als man nur will, wenn man nur die Seiten gross genug nehmen darf,
dennoch kann der Flächeninhalt eines Dreiecks, wie gross auch die Seiten genommen werden, nie
eine bestimmte Grenze überschreiten, ja sie nicht einmal erreichen. Alle meine Bemühungen,
einen Widerspruch, eine Inconsequenz in dieser Nicht-Euklidischen Geometrie zu finden, sind
fruchtlos gewesen, und das Einzige, was unserm Verstande darin widersteht, ist, dass es, wäre
sie wahr, im Raum eine \textls{an sich bestimmte} (obwohl uns unbekannte) Lineargrösse geben
müsste. Aber mir deucht, wir wissen, trotz der nichtssagenden Wort-Weisheit der Metaphysiker
eigentlich zu wenig oder gar nichts über das wahre Wesen des Raums, als dass wir etwas uns
unnatürlich vorkommendes mit \textls{Absolut Unmöglich} verwechseln dürfen. Wäre die
Nicht-Euklidische Geometrie die wahre, und jene Constante in einigem Verhältnisse zu solchen
Grössen, die im Bereich unserer Messungen auf der Erde oder am Himmel liegen, so liesse sie
sich a posteriori ausmitteln. Ich habe daher wohl zuweilen im Scherz den Wunsch geäussert,
dass die Euklidische Geometrie nicht die wahre wäre, weil wir dann ein absolutes Mass a priori
haben würden.

Von einem Manne, der sich mir als einen denkenden mathematischen Kopf gezeigt hat, fürchte ich
nicht, dass er das Vorstehende missverstehen werde: auf jeden Fall aber haben Sie es nur als
eine Privat-Mittheilung anzusehen, von der auf keine Weise ein öffentlicher oder zur
Öffentlichkeit füh%
